%% Table 1 — Data and Event Summary
%% Source: Official PJM Data Miner 2 + ERA5 CDS

\begin{table}[htbp]
\centering
\caption{Data summary and event definition. Training uses 2010--2013 PJM RTO hourly load with ERA5 great\_lakes\_core retrospective reanalysis weather. The 2014 test year includes the Jan~6--8 Polar Vortex cold-weather stress window and the Jun~16--18 summer near-peak window. Load values are in MW.}
\label{tab:data_summary}
\begin{tabular}{lcccccc}
\toprule
Window & Period & Load Mean & Load Max & Load Min & Temp Mean & HDH Mean \\
\midrule
Training & 2010--2013 & 91,034 & 158,043 & 49,558 & 52.3 & 12.7 \\
Test (full-year) & 2014 & 91,034 & 141,678 & 57,570 & 49.8 & 15.1 \\
Cold-event window & Jan 6--8, 2014 & 120,069 & 140,510 & 82,638 & 3.4 & 61.2 \\
Summer window & Jun 16--18, 2014 & 126,845 & 141,678 & 108,911 & 73.3 & 0.0 \\
\bottomrule
\end{tabular}
\end{table}

%% Table 2 — Model Summary
%% Source: Task07A + Task07B verified outputs

\begin{table}[htbp]
\centering
\caption{Models evaluated. All weather-informed models use ERA5 retrospective reanalysis (not operational forecast). For quantile models, post-hoc monotonic rearrangement was applied to correct quantile crossings (66\% of test hours affected).}
\label{tab:model_summary}
\begin{tabular}{lllll}
\toprule
Model & Type & Weather Input & Training & Role \\
\midrule
Persistence-1h & Naive baseline & None & None & Lower bound \\
Naive-24h & Seasonal naive & None & None & Seasonal benchmark \\
Naive-168h & Weekly naive & None & None & Weekly benchmark \\
GBoost & Point forecast & ERA5 reanalysis & 2010--2013 & Point benchmark \\
QR-GBT & Quantile (q01--q99) & ERA5 reanalysis & 2010--2013 & Prob. forecast \\
PJM Day-Ahead & Operational & PJM internal & N/A & External benchmark \\
\bottomrule
\end{tabular}
\end{table}

%% Table 3 — Point Forecast Metrics
%% Source: Task07A verified outputs

\begin{table}[htbp]
\centering
\caption{Point forecast performance. All values in MW. Jan~7 18:00 EPT is the official cold-event peak (140,510~MW). The GBoost and Linear models use ERA5 retrospective reanalysis weather features.}
\label{tab:point_metrics}
\begin{tabular}{lrrrrr}
\toprule
Model & Full-Year MAE & Full-Year RMSE & Vortex MAE & Vortex RMSE & Jan 7 18:00 Error \\
\midrule
Persistence-1h & 2,745 & 3,542 & 2,805 & 3,565 & $+$2,906 \\
Naive-24h & 5,960 & 8,105 & 15,598 & 18,645 & $+$9,973 \\
Naive-168h & 8,009 & 10,871 & 24,531 & 26,409 & $+$30,769 \\
Linear & 2,354 & 3,027 & 4,248 & 5,100 & $+$5,525 \\
GBoost & 721 & 994 & 1,705 & 2,037 & $-$1,388 \\
PJM Day-Ahead & 1,937 & 2,563 & 3,148 & 3,809 & $+$2,545 \\
\bottomrule
\end{tabular}
\end{table}

%% Table 4 — Probabilistic Forecast Metrics
%% Source: Task07B verified outputs

\begin{table}[htbp]
\centering
\caption{Probabilistic forecast performance of QR-GBT (q50 point, plus 90\% and 98\% prediction intervals). The 90\% PI coverage degrades from 86.8\% full-year to 66.7\% during the Polar Vortex. The official event peak (140,510~MW at Jan~7 18:00) falls outside the 98\% PI.}
\label{tab:prob_metrics}
\begin{tabular}{lrrrrr}
\toprule
Window & q50 MAE & Mean Pinball & 90\% PI Cov. & 98\% PI Cov. & Winkler (90\%) \\
\midrule
Full-year 2014 & 761 & 153 & 86.8\% & 97.3\% & 4,627 \\
Jan 6--8 Vortex & 2,130 & 473 & 66.7\% & 84.7\% & 15,476 \\
Jun 16--18 Summer & 1,060 & 209 & 86.1\% & 93.1\% & 6,355 \\
Top 1\% load hours & 2,539 & 502 & 63.2\% & 82.8\% & 16,168 \\
\bottomrule
\end{tabular}
\end{table}

%% Table 5 — Event Peak Prediction Check
%% Source: Task07B verified outputs

\begin{table}[htbp]
\centering
\caption{Prediction check at the official cold-event peak (Jan~7, 2014 18:00 EPT, 140,510~MW). Positive error means underprediction. The actual load exceeds the QR-GBT 98\% prediction interval.}
\label{tab:event_peak}
\begin{tabular}{lrrcc}
\toprule
Model & Prediction (MW) & Error (MW) & In 90\% PI? & In 98\% PI? \\
\midrule
PJM Day-Ahead & 137,965 & $+$2,545 & — & — \\
Persistence-1h & 137,604 & $+$2,906 & — & — \\
Naive-24h & 130,537 & $+$9,973 & — & — \\
GBoost point & 141,898 & $-$1,388 & — & — \\
QR-GBT q50 & 135,357 & $+$5,153 & Yes & No \\
QR-GBT q95 & 139,410 & $+$1,100 & No & No \\
QR-GBT q99 & 139,585 & $+$925 & No & No \\
\bottomrule
\end{tabular}
\end{table}
